% Manometro a mercurio a tubo aperto: il mercurio nel ramo aperto sta 19 cm più in
% alto. Alla stessa quota A e B hanno la stessa pressione (ud3-statica-manometro).
\begin{tikzpicture}[>=Stealth, font=\small]
  % recipiente con il gas
  \fill[lampo!20] (0,2.4) rectangle (1.4,3.8);
  \draw[thick] (0,2.4) rectangle (1.4,3.8);
  \node at (0.7,3.1) {gas, $p$};
  % mercurio nel tubo a U
  \fill[richiamo!45] (1.8,1.2) -- (1.8,0.75) arc[start angle=180, end angle=360, radius=0.75]
    -- (3.3,2.6) -- (3.0,2.6) -- (3.0,0.75) arc[start angle=360, end angle=180, radius=0.45]
    -- (2.1,1.2) -- cycle;
  % pareti
  \draw[thick] (1.4,2.9) -- (1.8,2.9) -- (1.8,0.75) arc[start angle=180, end angle=360, radius=0.75] -- (3.3,3.6);
  \draw[thick] (1.4,3.2) -- (2.1,3.2) -- (2.1,0.75) arc[start angle=180, end angle=360, radius=0.45] -- (3.0,3.6);
  % livelli
  \draw[dashed, richiamo] (1.8,1.2) -- (3.9,1.2);
  \draw[dashed, richiamo] (3.3,2.6) -- (3.9,2.6);
  \draw[|<->|] (3.75,1.2) -- (3.75,2.6) node[midway, right, align=left] {$\Delta h = \SI{19}{\centi\metre}$\\$= \SI{190}{\mmHg}$};
  \fill (1.95,1.2) circle (1.6pt);
  \node[left] at (1.75,1.35) {A};
  \fill (3.15,1.2) circle (1.6pt);
  \node at (3.52,1.4) {B};
  \draw[->, thick, attenzione] (3.15,4.3) -- (3.15,3.65) node[pos=0, above] {$p_\text{atm} = \SI{760}{\mmHg}$};
  % pressioni in A e B
  \node[below, align=center] at (2.55,-0.1) {$p_\text{A} = p$ \qquad $p_\text{B} = p_\text{atm} + \rho g\,\Delta h$};
\end{tikzpicture}
